\section{Discussion and limitations}\label{sec:discussion}

\paragraph{What the theorems cover.}
The results of \cref{sec:calculus,sec:identity,sec:float,sec:verification,sec:case} are about
the specification and a core calculus, proven by hand. The Rust implementation is not verified;
its agreement with the laws is tested (\cref{tab:g1}), and an independent review found that it
violates two of our side conditions (\cref{sec:status}). The calculus omits recursion, sum
types, generics and explicit memory regions, which the MVP rejects, and treats each operation
through its typing, domain and meaning without restating the opcode table.

\paragraph{Assumptions.}
\Cref{thm:hash-char,thm:hash-meaning,lem:eff-never-merged} rest on collision freeness, which
holds except with the probability of \cref{lem:birthday} in the random-oracle idealization, a
heuristic for BLAKE3. \Cref{thm:soundness} rests on a faithful encoding of each operation
(\cref{ass:faithful}), on the solvers, and, for compiled code, on the code generator, LLVM, the
runtime and the operating system; the specification mitigates these with two-solver agreement
in release builds, replay of every counterexample, mutation testing of the verifier and
differential testing of every build, none of which has run yet. Independently checkable
certificates are planned.

\paragraph{Two corrections to the specification.}
Formalizing the specification exposed two imprecisions that matter for soundness, and we
recommend recording both as decisions before the verifier is built. First, an obligation's
context must contain only the facts of operations that precede it in a topological order
(\cref{def:ordered}); the current wording allows a program that traps to verify
(\cref{prop:ordering-needed}). Second, the facts of a \kw{spec:} application must be guarded by
the specification's precondition (\cref{prop:foldspec}). Four smaller points are matters of
precision: the cost expression is an upper bound, exact only without guards
(\cref{prop:cost}); the statement that canonical forms coincide exactly for texts related by the
listed moves needs the further identifications stated after \cref{cor:canon}; the
well-definedness condition of a fold term must require $0\le k$ as well as
$k\le\mathrm{len}(xs)$; and a kernel applied inside a formula can read an unspecified value
without a well-definedness obligation. The last two are not unsound, but they leave such a
contract without a single value.

\paragraph{Rules planned after the MVP.}
The specification already lists proof rules for \op{SCAN\_PREFIX} and \op{MATCH\_DISJOINT}.
As written, the first requires only associativity of the operator and the second only that the
arms are exhaustive and disjoint; neither requires the operator's or the arm's precondition to
hold where it is applied. For a partial operator such as checked addition, a prefix scan can
trap. These operations are not part of the MVP, but their rules need precondition goals, in the
style of \cref{tab:rules}, before they are implemented.

\paragraph{The contract--intent gap.}
A verified node satisfies its contract, not its author's intent. The case study shows both
sides: the sign clause catches a clamp that maps NaN to $-1$, which a range check alone would
accept, yet the same clause cannot tell $-0$ from $+0$ (\cref{rem:contract}). The specification
counts this among its gravest risks, renders every contract in plain words for human
confirmation, and leaves its measurement to a research track.

\paragraph{Floating point beyond the MVP.}
The MVP has no transcendental functions, because the solvers' floating-point theory has none.
They are planned as uninterpreted functions axiomatized from a verified, correctly rounded math
library, followed by a final clamp. Whether float-heavy obligations solve in acceptable time is
an empirical question that the verifier will answer; the specification falls back to an explicit
guard, counted as residual runtime cost, when a solver returns unknown.

\paragraph{What has not been measured.}
We make no claim about solver times, the share of obligations proven automatically, the speed of
generated code, the overhead of residual guards, token savings of the textual form, or the rate
at which models write verified nodes. The project's gates G3 to G8 define how each will be
measured; none has run.

\section{Conclusion}\label{sec:conclusion}

\VG{} treats a program as a content-addressed, contract-carrying dataflow graph that an AI agent
writes, checks and repairs through structured judgments. We have given the design formal
foundations: a core calculus; a characterization of when two bodies share a hash, with the
capacity condition and the linearity of effect tokens as the precise hypotheses; a proof that no
efficient canonical form can capture semantic equivalence, which justifies spec classes; a
one-NaN semantics under which IEEE operations are congruences and the commutative ones commute
bit for bit; soundness of the combinator rules and of the exact integer encoding, with an
ordering proviso that we showed necessary; and a complete proof of the case study's kernel,
confirmed by an exhaustive enumeration. The prototype implements the core representation, and an
independent review's findings fall exactly on our side conditions. The verifier, the compiler
and the agent loop remain to be built, and the measurements that would justify claims about them
remain to be made.
